\section{The Erdős--Straus workbench}\label{sec:es}

The Erdős--Straus conjecture asserts that $4/n=1/x+1/y+1/z$ has a solution in positive integers for every $n\ge2$; it suffices to treat primes $n=p$.
The workbench \texttt{erdos-straus-foundation} states that it claims neither a proof nor a counterexample.
Its collective author name is The Clankers; the mod-107 observation behind its result R06 is credited to u/CommonCareful3149, further contributions to u/UmbrellaCorp\_HR, and the diagrams behind R05 to u/CivQ17 (README, ``Attribution and collaboration''; \texttt{results/records.json}). The workbench's fuller contribution record (\texttt{CONTRIBUTIONS.md}) also credits u/Far-Lifeguard519 for the initiating computational investigation that prompted the project, u/TextBackground496 and u/JGPTech for source ideas, and the owner (KokunoYumeto, u/lepthymo) for proposing the research directions and examples, directing and developing the investigations with AI, and specifying the proof and source-fidelity requirements.
The results of its 22 September release are described as written proofs with Python replays, without a Lean build or independent human review.
Audit note: \note{43\_}; first-pass report \note{43a\_}.

Throughout this section $p$ is an odd prime, and a \emph{solution} is a triple of positive integers with $4/p=1/x+1/y+1/z$. The workbench works with $p=12h+1$ ($h\ge1$).

\subsection{The shell criterion}

For an integer $a>p/4$ not divisible by $p$ put $R_a=4a-p$, $S_a=pa$ and
\[
E_a=\{u: u\mid a^2,\ R_a\mid 4u+1\},\qquad M_a=\{u: u\mid a^2,\ R_a\mid u+a\}.
\]
A \emph{shell} $a$ is \emph{occupied} if $E_a\cup M_a\neq\emptyset$.

\begin{theorem}[shell criterion]\label{thm:shell}
\begin{enumerate}[label=(\alph*)]
\item If $x\le y\le z$ is a solution, then $p/4<x\le3p/4$. More precisely, $x<p/2$ and $x<y$, unless $p\equiv3\pmod4$ and $(x,y,z)=\bigl(\frac{p+1}2,\frac{p+1}2,\frac{p(p+1)}4\bigr)$. For $p=12h+1$ this gives $3h+1\le x\le 6h$.
\item For each integer $a>p/4$ not divisible by $p$ (in particular each $a$ with $p/4<a<p$), the ordered pairs $(y,z)$ of positive integers with $1/a+1/y+1/z=4/p$ number exactly $2|E_a|+|M_a|$. A pair with $y=z$ occurs if and only if $R_a=1$ (possible only for $p\equiv3\pmod4$ and $a=(p+1)/4$); otherwise $|M_a|$ is even and the unordered pairs number $|E_a|+\tfrac12|M_a|$. For $p\equiv1\pmod4$ this is always the case.
\item Hence the equation is solvable at $p$ if and only if some shell $a$ with $p/4<a<p/2$ is occupied.
\end{enumerate}
\end{theorem}
\status{Generalised here: the workbench's ES-01/ES-02 (\texttt{bounded\_transport.tex}), audited, state it for $p=12h+1$ and $a\in[3h+1,9h]$; the proof below gives every odd prime, every $a>p/4$ with $p\nmid a$, and the halved range $x<p/2$ ($a\in[3h+1,6h]$ for $p=12h+1$). The halved range at $p\equiv1\pmod{12}$, with $x$ unique, is already a reduction step of the workbench's item~22, checked by reading in the first pass (\note{43a\_} §2.22), and items 6 and 7 work in it; the workbench's later packages attribute these coordinates to Elsholtz and Tao \cite{ElsholtzTao}, §2. Checked against a brute-force count of the ordered pairs on all $4{,}726$ pairs $(p,a)$ with $p<260$ an odd prime and $p/4<a<p$ (a pair with $y=z$ exactly in the $29$ cases with $R_a=1$) and on all $4{,}303$ pairs with $p<120$, $p/4<a\le3p$ and $p\nmid a$; on all $4{,}398$ shells of the primes \mbox{$p\equiv1\pmod{12}$} below $700$, and in the first pass on all $8{,}124$ shells below $1000$; (a) against all $2{,}259$ solutions at the odd primes below $400$; every prime \mbox{$p\equiv1\pmod{12}$} below $30{,}000$ has an occupied shell.}

\begin{proof}
(a) The other two terms are positive, so $1/x<4/p$; and $x$ is the least denominator, so $4/p\le3/x$. Hence $p/4<x\le3p/4$.
Suppose $x>p/2$, so $x\ge(p+1)/2$ as $p$ is odd, and $x<p$. Then $2/y\ge1/y+1/z=4/p-1/x>2/p$, so $y<p$. From $4xyz=p(xy+yz+zx)$ the prime $p$ divides $xyz$, and as $0<x\le y<p$ it divides $z=pz'$; then $xy(4z'-1)=pz'(x+y)$ with $p\nmid xy$, so $4z'-1=pk$ with $k\ge1$. Put $N=pk+1=4z'$; then $1/x+1/y=4k/N$, that is, $(4kx-N)(4ky-N)=N^2$ with both factors positive. But $4kx-N\ge2k(p+1)-N=N+2(k-1)$, and likewise for $y$; so $k=1$ and $x=y=N/2=(p+1)/2$, with $z=p(p+1)/4$, which needs $p\equiv3\pmod4$. Finally, if $x<p/2$ and $y=x$, then $1/x+1/y>4/p$, which is impossible. For $p=12h+1$ the integers in $(p/4,p/2)$ are $3h+1,\dots,6h$.

(b) Write $R=R_a$, $S=S_a$. Then $R\ge1$ is odd, because $p$ is. Since $p\nmid a$, $\gcd(a,R)=\gcd(a,p)=1$ and $\gcd(p,R)=\gcd(p,4a)=1$, so $S$ is a unit modulo $R$.
The equation $1/y+1/z=R/S$ is equivalent to $(Ry-S)(Rz-S)=S^2$, and both factors are positive because $1/y$ and $1/z$ are smaller than $R/S$. So the ordered pairs correspond bijectively to ordered factorisations $S^2=dd'$ into positive divisors with $d\equiv d'\equiv-S\pmod R$, through $y=(d+S)/R$, $z=(d'+S)/R$. If $d\mid S^2$ and $d\equiv-S$, then $d$ is a unit and $d'=S^2/d\equiv S^2(-S)^{-1}=-S$; so the pairs are counted by $\#\{d\mid S^2: d\equiv-S\pmod R\}$.
Since $\gcd(p,a)=1$, each divisor of $S^2=p^2a^2$ is uniquely $d=p^iv$ with $i\in\{0,1,2\}$ and $v\mid a^2$. Using $4a\equiv p\pmod R$:
for $i=2$, $p^2v\equiv-pa\iff pv\equiv-a\iff 4pv\equiv-p\iff R\mid 4v+1$, so these $v$ form $E_a$;
for $i=1$, $pv\equiv-pa\iff R\mid v+a$, so these $v$ form $M_a$;
and $d\mapsto S^2/d$ maps the layer $i=2$ bijectively onto the layer $i=0$ while preserving the congruence. The total is $2|E_a|+|M_a|$.
The swap $(y,z)\mapsto(z,y)$ is $d\mapsto S^2/d$; it maps each layer $i$ to the layer $2-i$, so a fixed point lies in the layer $i=1$, where it is $v\mapsto a^2/v$ with only possible fixed point $v=a$. That needs $R\mid 2a$, i.e. $R\mid2$ (as $\gcd(R,a)=1$), i.e. $R=1$. If $R\ne1$ the swap has no fixed point, $|M_a|$ is even, $y\ne z$, and the unordered pairs number $|E_a|+\tfrac12|M_a|$. For $p\equiv1\pmod4$, $R\equiv-p\equiv3\pmod4$, so $R\ne1$.

(c) For $p\equiv1\pmod4$ it follows from (a) and (b). For $p\equiv3\pmod4$ both sides hold: the shell $a=(p+1)/4$ lies in $(p/4,p/2)$ and has $R_a=1$, so $u=a$ lies in $M_a$.
\end{proof}

The workbench also gives the exact bookkeeping around this criterion for $p=12h+1$: an explicit bijection between a tagged set and the ordered solutions, its inverse read off from the $p$-adic valuation of $Ry-S$, and the global statement that a counterexample prime would satisfy the finite, computable condition $E_a=M_a=\emptyset$ for every $a$ (\note{43a\_} §1.2; re-derived there).

\subsection{The first two shells}

\begin{proposition}[the residual-3 shell]\label{prop:r3}
Let $a=3h+1=(p+3)/4$, so $R_a=3$. Let $P_1=\prod(2v_\ell(a)+1)$ over the primes $\ell\equiv1\pmod3$ dividing $a$, and $P_2$ the same product over the primes $\ell\equiv2\pmod3$. Then $E_a=M_a=\{u\mid a^2: u\equiv2\pmod3\}$ and $|E_a|=|M_a|=P_1(P_2-1)/2$. So a solution with least denominator $a$ exists if and only if $a$ has a prime factor $\equiv2\pmod3$, and there are then exactly $3P_1(P_2-1)/4$ unordered solutions with least denominator $a$.
\end{proposition}
\status{Proved here; the workbench's residual-3 sieve (\texttt{first\_two\_shell\_sieve.tex}), audited. Checked on all 4{,}472 primes $p\equiv1\pmod{12}$ below $2\cdot10^5$, and in the first pass on all 19{,}564 below $10^6$.}

\begin{proof}
Since $a\equiv1\pmod3$, $3\nmid a$. Modulo $3$ both $4u+1$ and $u+a$ are $\equiv u+1$, so both conditions say $u\equiv2\pmod3$. A divisor $u=\prod\ell^{e_\ell}$ of $a^2$ has $u\equiv(-1)^{\sum e_\ell}\pmod3$, the sum over the primes $\ell\equiv2\pmod3$. Their exponent vectors ($0\le e_\ell\le2v_\ell(a)$) number $P_2$, and the difference between the numbers with even and with odd sum is $\prod_\ell\sum_{e=0}^{2v_\ell}(-1)^e=1$; so $(P_2-1)/2$ have odd sum, and the other primes contribute the factor $P_1$. By Theorem~\ref{thm:shell}(b) the unordered solutions containing $a$ number $\tfrac32P_1(P_2-1)/2$, and $a=3h+1$ is the least denominator of each of them because it is the smallest possible one. Finally $P_2>1$ exactly when $a$ has a prime factor $\equiv2\pmod3$. (The count is an integer: $a\equiv1\pmod 3$ makes the total exponent of those primes even, so $P_2\equiv1\pmod4$.)
\end{proof}

\begin{proposition}[the residual-7 shell]\label{prop:r7}
Let $a=3h+2=(p+7)/4$, so $R_a=7$ and $7\nmid a$. Let $n_3$ be the total exponent in $a$ of the primes $\equiv3\pmod7$ and $n_{24}$ that of the primes $\equiv2$ or $4\pmod 7$. The shell is occupied if and only if $a$ has a prime factor $\equiv5$ or $6\pmod7$, or $n_3\ge3$, or $n_3\ge1$ and $n_{24}\ge1$.
\end{proposition}
\status{Proved here; the workbench's residual-7 sieve, audited, together with its count as coefficients of $\prod_\ell\sum_{e=0}^{2v_\ell(a)}T^{\lambda(\ell)e}$ in $\ZZ[T]/(T^6-1)$. Checked on all 4{,}472 primes $p\equiv1\pmod{12}$ below $2\cdot10^5$ (and below $10^6$ in the first pass).}

\begin{proof}
The exterior condition $7\mid4u+1$ is $u\equiv-4^{-1}\equiv5=3^5$, and the middle condition is $u/a\equiv-1=3^3\pmod7$; in base $3$ the discrete logarithms of $1,\dots,6$ are $0,2,1,4,5,3$.
\emph{Sufficiency.} A prime $\ell\equiv5$ dividing $a$ gives $u=\ell$; a prime $\ell\equiv6$ gives $u=a\ell$. If $n_3\ge3$, a product of primes $\equiv3$ dividing $a^2$ with total exponent $5$ gives $u\equiv3^5$. If primes $t\equiv3$ and $\ell\in\{2,4\}\pmod 7$ divide $a$, then $u=a\ell t$ (when $\ell\equiv2=3^2$) or $u=a\ell/t$ (when $\ell\equiv4=3^4$) divides $a^2$ and has $u/a\equiv3^3$.
\emph{Necessity.} If none of the conditions holds, every prime factor of $a$ is $\equiv1,2,3$ or $4$, and either $n_3=0$ or ($1\le n_3\le2$ and $n_{24}=0$). If $n_3=0$, $a$ and all divisors of $a^2$ lie in the subgroup $\{1,2,4\}$ of squares, which contains neither $5$ nor $-1$. Otherwise $a\equiv3^{n_3}$ and a divisor of $a^2$ is $\equiv3^e$ with $0\le e\le2n_3\le4$; the exterior target needs $e\equiv5\pmod6$ and the middle one $e-n_3\equiv3\pmod6$ with $|e-n_3|\le2$, both impossible.
\end{proof}

The shells with $R=3$ and $R=7$ exist exactly for the primes $p\equiv1\pmod4$, with $a=(p+3)/4$ and $a=(p+7)/4$ (\textsf{Generalised here}; checked on all $1{,}611$ such primes below $3\cdot10^4$). Proposition~\ref{prop:r7} holds for all of them, since its proof never uses $p\equiv1\pmod3$. Proposition~\ref{prop:r3} holds as stated for $p\equiv1\pmod{12}$ only. For $p\equiv5\pmod{12}$, $a=(p+3)/4\equiv2\pmod3$, so modulo $3$ the middle condition reads $u\equiv1$; then $E_a=\{u\mid a^2:u\equiv2\}$ and $M_a=\{u\mid a^2:u\equiv1\}$, with $|E_a|=P_1(P_2-1)/2$ and $|M_a|=P_1(P_2+1)/2$ by the parity count of the proof above, and there are exactly $|E_a|+\frac12|M_a|=P_1(3P_2-1)/4$ unordered solutions with least denominator $a$ (an integer: the total exponent in $a$ of the primes $\equiv2\pmod3$ is odd, so $P_2\equiv3\pmod4$). That shell is always occupied ($1\in M_a$), and its occupancy criterion holds trivially, since $a\equiv2\pmod3$ has a prime factor $\equiv2\pmod3$; those primes are solvable anyway, since $p\equiv2\pmod3$.

\begin{proposition}[the two-shell survivors]\label{prop:survivors}
Let $p\equiv1\pmod{12}$ be prime. Both shells $a=(p+3)/4$ and $a=(p+7)/4$ are unoccupied if and only if $p\equiv1\pmod{24}$, every prime factor of $(p+3)/4$ is $\equiv1\pmod3$, and every prime factor of $(p+7)/4$ is $\equiv1,2$ or $4\pmod7$. Every such $p$ is a square modulo $7$, so $p\equiv1$, $25$ or $121\pmod{168}$.
\end{proposition}
\status{Proved here; the first pass stated the congruence mod $24$ as an observation on its computation. Checked on all primes $p\equiv1\pmod{12}$ below $2\cdot10^5$.}
\begin{proof}
If $p\equiv13\pmod{24}$, then $(p+3)/4$ is even and $2\equiv2\pmod3$, so the first shell is occupied (Proposition~\ref{prop:r3}). If $p\equiv1\pmod{24}$, the first shell is unoccupied exactly when every prime factor of $(p+3)/4$ is $\equiv1\pmod3$; and $(p+7)/4$ is even, so $n_{24}\ge1$ and Proposition~\ref{prop:r7} says that the second shell is occupied exactly when $(p+7)/4$ has a prime factor $\equiv3$, $5$ or $6\pmod7$, a non-residue. If it has none, $(p+7)/4$ is a square modulo $7$, and so is $p\equiv4\cdot\frac{p+7}4$. Among the classes $\equiv1\pmod{24}$ modulo $168$, the squares modulo $7$ are $1$, $25$ and $121$.
\end{proof}
Of the primes $p\equiv1\pmod{12}$ below $10^6$, exactly $989$ have both shells unoccupied, the first being $1129$, $1201$, $2521$ (first pass).

\subsection{Which primes remain}

The workbench proves solvability for $p\equiv2\pmod 3$, $p\equiv3\pmod4$ and $p\equiv13\pmod{24}$ (scope identities that it does not advertise as new) by the families $(p,\frac{p+1}3,\frac{p(p+1)}3)$, $(\frac{p+1}4,\frac{p(p+1)}2,\frac{p(p+1)}2)$ and $(a,\frac{pa}2,pa)$ with $a=\frac{p+3}4$ (even when $p\equiv13\pmod{24}$); so its own reduction leaves the primes $p\equiv1\pmod{24}$. Its later packages restrict to the \emph{hard} primes $p\equiv1,121,169,289,361,529\pmod{840}$, which are exactly the unit squares modulo $840$ (a group of order $1\cdot1\cdot2\cdot3=6$). The reduction to them is classical (Mordell \cite{Mordell}) and is not proved in the texts read; five identities prove it.

\begin{proposition}[the reduction modulo $840$]\label{prop:mod840}
Let $p\equiv1\pmod{24}$ be prime, and let $(A,B,R)$ be $(1,2,7)$, $(2,1,7)$, $(1,1,7)$, $(1,2,15)$ or $(2,1,15)$ according as $p\equiv3$, $5$, $6\pmod7$ or $p\equiv2$, $3\pmod5$. Then $D=(p+R)/(4AB)$ and $C=(A+pB)/R$ are positive integers and
\[
\frac4p=\frac1{ABD}+\frac1{ACD}+\frac1{pBCD}.
\]
So the equation is solvable at every prime $p\equiv1\pmod{24}$ that is not a square modulo $840$.
\end{proposition}
\status{Proved here (the reduction is classical); the first pass covered the $18$ non-square unit classes $\equiv1\pmod{24}$ by the same kind of families. Checked on all primes $p\equiv1\pmod{24}$ below $3\cdot10^5$.}
\begin{proof}
With $RC=A+pB$ and $4ABD=p+R$: $\frac1{ABD}+\frac1{ACD}+\frac1{pBCD}=\frac{pC+pB+A}{pABCD}=\frac{C(p+R)}{pABCD}=\frac{p+R}{pABD}=\frac4p$. Since $p\equiv1\pmod8$, $8$ divides $p+7$ and $p+15$, so $4AB\mid p+R$; $R\mid A+pB$ is the stated congruence modulo $7$ or $5$, together with $A+pB\equiv A+B\equiv0\pmod3$ when $R=15$. The primes $p\equiv1\pmod{24}$ left over have $p\equiv1,2,4\pmod7$ and $p\equiv\pm1\pmod5$, which are exactly the unit squares modulo $840$.
\end{proof}

\subsection{The structural items}

The workbench's 28 structural items (17--23 September) are listed in \note{43a\_} §2. Those audited:
\begin{itemize}
\item \emph{Item 6} (audited). With $L_a=8\tau(a)^2/\varphi(R_a)-\mathcal E_a$, where $\mathcal E_a$ is the residue-collision energy of the $2\tau(a)$ divisors of $pa$ modulo $R_a$: for $p\equiv1\pmod4$, $p\ge2^{20}$, $\sum_{a\in(p/4,p/2)}L_a\le-\tfrac1{24}p\log p$. So the lower bound that keeps only the principal character, summed over the shells, is negative for large $p$: that method cannot prove occupancy.
\item \emph{Item 7} (audited). For $p\equiv1\pmod 8$, bounds in terms of the least quadratic nonresidue; at the hard primes a middle state $(a,u)$ on a shell $p/4<a<p/2$ has $a\le3(p+3)/11$, and an exterior state satisfies $22(p+1)\le23(p-2a+1)^2$. The workbench says this localises solutions and does not prove the region occupied.
\item \emph{Item 14} (audited). For $h'=4j-1>t$ the divisors $W\mid j^2$ with $W\equiv t\pmod{h'}$ are canonical, companion or finite templates; the finite ones satisfy the sharp bound $h'\le t^2-3t+1$, with equality exactly in one parametrised family, attained at infinitely many primes $p\equiv1\pmod{840}$. For every $t\ge4$ infinitely many hard primes have an exterior state whose same-grade middle box is empty.
\item \emph{Item 21} (audited). For a proper rational trace source the complete set of same-word residual returns is $\{k\mid R: d\mid k,\ k\equiv1\pmod{4K(u)}\}$; a uniform return exists exactly for the nine words $u\mid36$.
\item \emph{Item 22} (audited, certificate). For every solution at $p\equiv1\pmod{12}$ an explicit symmetric polynomial $\mathfrak N(p,x,y,z)$ of degree $8$ satisfies $\mathfrak N<-\frac{125873811}{262144}p^8$ ($\approx-480.17\,p^8$), and $\mathfrak N<-6960.19\,p^8$ when $p\equiv1\pmod{24}$; the proof reduces to finite polynomial positivity certificates, recomputed independently. It is a statement about existing solutions, not about their existence.
\item \emph{Item 26} (audited in part: the least example and the deduction from it; its prime-square trace theorem at $p\equiv1\pmod{24}$ and its terminality statement read and correct as read, with the identity $16(4q^3+1)-(4q+1)(16q^2-4q+1)=15$ checked; the factor-sum, Pell and fixed-tail-sum sections not checked). $p=67369$ is the least prime $p\equiv1\pmod{24}$ whose first ``trace shell'' ($a=16849$, $R=27$) precedes its first full shell ($a=16850$, $R=31$); so no universal map from rational traces to solutions avoids increasing $a$.
\item \emph{Item 4} (conditional on a computer enumeration of $1{,}381{,}117{,}764$ rooted profiles, $5{,}922{,}259$ of them range-compatible, run by two implementations of the workbench and not re-run here). At every hard prime, from any nonresidue prime, factoring at most 25 numbers $(p+\sigma qt^2)/4$ with $t\in\{1,3,5,7,9\}$ yields at least six distinct nonresidue primes. A direct test of all 108 hard primes below $40{,}000$ agrees.
\end{itemize}
Items 2, 3 and 8 are audited (certificate); the $22$-digit prime $7510085481569082811681$ of item~3 is proved prime here by a Pocklington chain ($p-1=2^5\cdot3\cdot5\cdot q_1$, $q_1-1=2\cdot3\cdot5\cdot351707\cdot q_2$, $q_2-1=2\cdot3\cdot17^2\cdot5051\cdot169307$, witnesses below $12$ at every level, the leaves by trial division). Item 23 is audited in part (its certificates; the deduction read in part); items 11--13, 15 and 20 use imported material (Section~\ref{sec:intersections}); the remaining 17 items, including these five, are located. The results bench R01--R10 and the arguments X1--X7 are checked in \note{43a\_} §3; one source passage error (X5: $289^3\equiv169$, not $-1$, modulo $840$) is recorded by the workbench itself.

\subsection{The \texorpdfstring{$(3,4,\infty)$}{(3,4,infinity)} monodromy covers}\label{ssec:es-s6}

A package of the workbench (ES-09) imports three integer matrices $A_1,A_2,A_\infty$ with $A_1^3=A_2^4=A_1A_2A_\infty=I$ from the $S^6$ manuscript \cite{AlpogeS6} and nothing else from it.
For odd levels $D$, which is the source's standing hypothesis, the induced affine action on $(\ZZ/D)^3$ is transitive when $3\nmid D$ and has two orbits, of sizes $D^3/9$ and $8D^3/9$, when $3\mid D$; for $3\nmid D$ the cover is connected of genus $(5D^3-12D^2-17D+24)/24$, and for $3\mid D$ the two orbits give components of genera $(5D^3-12D^2-81D+216)/216$ and $(5D^3-12D^2-9D+27)/27$.
These statements are false for even $D$ (at $D=2$ the orbits have sizes $2$ and $6$; found by the referee of the notes), which the ES application never uses, since there $D$ divides the odd number $R_a$.
The package also shows that the equation fails at $p$ if and only if a denominator-labelled cover has no section. That equivalence is correct, and it is Theorem~\ref{thm:shell} again: a section exists exactly when the tail denominator $R_a/\gcd(R_a,4u+1)$ or $R_a/\gcd(R_a,u+a)$ equals $1$.
\status{Audited: the proof read; the relations, the orbits and the genera also recomputed from cycle counts for $D=1,3,\dots,15$. The equivalence is a re-encoding of Theorem~\ref{thm:shell}; the assessment is this audit's (the package presents it as a complete failure object).}

\subsection{Negative results, with exact scope}

Each row says which proposed implication fails and what the failure does not show. Every certificate prime in the table is itself solvable.

\begin{center}\small
\begin{longtable}{@{}>{\raggedright\arraybackslash}p{0.19\linewidth}>{\raggedright\arraybackslash}p{0.29\linewidth}>{\raggedright\arraybackslash}p{0.21\linewidth}>{\raggedright\arraybackslash}p{0.21\linewidth}@{}}
\toprule
Source & What fails & Certificate & What it does not show\\
\midrule
\endhead
bounded transport & every occupied shell has a solution with unary layer $A\le2$ & $p=37$, $a=12$: only $u=8$ & the global cutoff ($a=10$ works at $37$)\\
first-two-shell sieve & occupancy of the $R=7$ shell gives a first-two-layer hit & $p=373$, $a=95$ & anything about other shells\\
item 2 & every canonical closed component contains a local hit & $p=2521$: a 7-cycle sink & $2521$ is solvable\\
item 3 & strict canonical descent ends at a hit & a closed triangle at a 22-digit prime & that prime is solvable\\
item 6 & the principal-only bound, summed over shells, is positive for large $p$ & $\sum L_a\le-\frac1{24}p\log p$ & anything about the true counts\\
item 6 & a nonnegative reweighting of it certifies occupancy & $p=87481$: all $21{,}870$ shells have $L_a\le0$ & only this certificate family fails\\
item 14 & same-grade middle repair from exterior states with $\alpha=-4t$, $t\ge4$ & $p=7840561$ & those primes are solvable\\
item 21 & a same-word residual return for words $u\nmid36$ & $3{,}041$ of the $3{,}654$ sources with $u\nmid36$ at primes below $2600$ & endpoint solutions exist\\
item 26 & a trace-to-solution map never increases $a$ & $p=67369$ & $67369$ is solvable\\
\bottomrule
\end{longtable}
\end{center}
All of these are statements about methods or certificate families. None is a counterexample to the conjecture, and none bears on whether it is true.
